\chapter{Antibody Serology}

\section*{What Is This Test?}
Antibody serology is the detection and measurement of antibodies (immunoglobulins) in serum or other body fluids as markers of past exposure, current or recent infection, immune status, or autoimmune disease. The humoral immune response to infection or vaccination produces antibodies that can be detected and characterized by their target antigen specificity, immunoglobulin class (IgM, IgG, IgA, IgE), concentration (titer/ index/ IU/mL), and functional activity (neutralization, opsonization, complement fixation). In infectious disease serology, the presence of specific IgM antibodies suggests recent or acute infection (IgM appears first, within days to 1--2 weeks of infection), while IgG antibodies indicate past infection or vaccination and typically persist for years to life. IgG seroconversion (negative to positive) or a fourfold or greater rise in antibody titer between acute and convalescent serum samples (collected 2--4 weeks apart) is diagnostic of recent infection. Serology is used for: screening for immunity (VZV, measles, mumps, rubella, HAV, HBV anti-HBs), diagnosis of infections that are difficult to culture or detect by NAAT (syphilis, Lyme disease, EBV, CMV, HIV---with the 4th-generation assay also detecting p24 antigen), diagnosis of autoimmune diseases (ANA, rheumatoid factor, ANCA), and epidemiological surveillance. Serologic testing methods include enzyme immunoassay (EIA/ELISA), chemiluminescent immunoassay (CLIA), immunofluorescence assay (IFA), agglutination assays (TPPA for syphilis, RPR, Weil-Felix---historical), immunoblot (Western blot, line blot), and neutralization assays (PRNT for flaviviruses, HAI for influenza).

\section*{Alternative Names}
Serology; Antibody Testing; Immunoassay; Serologic Testing; IgG/IgM Panel; Immune Status Testing; Antibody Titers; Antibody Index

\section*{Principle}
The fundamental principle of antibody serology is the specific, high-affinity binding of an antibody (in the patient's serum) to its cognate antigen (viral protein, bacterial protein, autoantigen, etc.), with detection of this binding by one of several methods. Enzyme-linked immunosorbent assay (ELISA/EIA): Antigen is immobilized on a solid phase (microtiter well or microparticle). Patient serum is added; if specific antibodies are present, they bind to the antigen. After washing, an enzyme-conjugated anti-human immunoglobulin (anti-IgG, anti-IgM) detection antibody is added, which binds to the patient's antibody. A chromogenic substrate is added; the enzyme converts it to a colored product, and the optical density (absorbance) is measured. The signal is proportional to the amount of antibody present. Results are reported as a signal-to-cutoff (S/CO) ratio or index value; a ratio above the cutoff (typically 1.0 or 1.1) is reactive/positive. CLIA: Similar to ELISA but uses a chemiluminescent substrate that emits light; measured by a luminometer. More sensitive and has a wider dynamic range than colorimetric ELISA. Used in most high-throughput automated platforms (Roche cobas, Abbott Architect, Siemens Atellica, Ortho Vitros). Immunofluorescence assay (IFA): Antigen-expressing cells or tissue sections are fixed on a glass slide. Patient serum is applied and incubated. After washing, fluorescein-labeled anti-human immunoglobulin is added and the slide is examined under a fluorescence microscope. The pattern and intensity of fluorescence are reported (e.g., ANA: homogeneous, speckled, nucleolar, centromere patterns). IFA is labor-intensive and subjective but remains the gold standard for many autoantibody tests. Agglutination: Particles (gelatin, latex, charcoal) are coated with antigen. When mixed with serum containing specific antibodies, the particles agglutinate (clump), producing visible aggregation. Examples: TPPA (\textit{T. pallidum} particle agglutination for syphilis), RPR (charcoal-cardiolipin particles for syphilis), Weil-Felix (Proteus antigens for rickettsial disease---historical). Rapid/Point-of-care immunochromatographic tests: Lateral flow assays in which patient serum migrates along a strip containing immobilized antigen and a detection antibody-colloidal gold conjugate; a visible line appears if antibodies are present (e.g., HIV rapid test, Monospot). Immunoblot: Antigens are separated by electrophoresis (by molecular weight), transferred to a membrane, and incubated with patient serum. Antibodies binding to specific protein bands are detected by an enzyme-conjugated anti-human Ig with a colorimetric substrate. Examples: Western blot (historical for HIV and Lyme), line blot (ANA profile, ANCA profile, myositis panel, paraneoplastic panel). Neutralization assays: The ability of patient antibodies to neutralize the biological activity of the pathogen in cell culture or animal models. Examples: PRNT (plaque reduction neutralization test for flaviviruses---Zika, dengue, WNV), HAI (hemagglutination inhibition for influenza), toxin neutralization (ASO). Neutralization assays measure functional, protective antibodies and are the gold standard for immunity, but are labor-intensive, require live virus culture (BSL-3 for some agents), and are performed only at reference laboratories.

\section*{History}
The concept of humoral immunity and antibodies dates to Emil von Behring and Kitasato Shibasaburō (1890), who discovered antitoxins against diphtheria and tetanus. The Wassermann test for syphilis (complement fixation, 1906) was the first widely used serologic test for an infectious disease. The ELISA was invented by Peter Perlmann and Eva Engvall in 1971 \cite{engvall1971elisa}, revolutionizing serology by replacing radioimmunoassay (RIA) with a safer, non-radioactive method. Automated high-throughput immunoassay platforms were introduced in the 1990s--2000s. The HIV Western blot (1987--2014) was the standard confirmatory test for 27 years before being replaced by the differentiation immunoassay + NAT algorithm. Recombinant antigen technology (1980s--present) improved specificity by replacing crude whole-cell or tissue-extracted antigens with purified, well-defined recombinant proteins. The application of serology to COVID-19 (2020) for epidemiologic surveys, vaccine response assessment, and diagnosis of past infection demonstrated both the power and the limitations of serology (slow to become positive, variable durability, and cross-reactivity).

\section*{Why This Test Is Ordered}
Serology is ordered for a vast array of clinical scenarios detailed in the respective disease- and system-specific chapters:
\begin{itemize}
  \item \textbf{Immune status screening:} VZV IgG, measles IgG, mumps IgG, rubella IgG (prenatal panel and pre-employment healthcare worker screening), HAV IgG (pre-vaccination or post-exposure assessment), HBV anti-HBs (post-vaccination titer at 1--2 months to confirm immunity: $\geq$10 mIU/mL is protective), HBV anti-HBc (to distinguish vaccine immunity from infection-acquired immunity)
  \item \textbf{Diagnosis of acute infection:} IgM class antibodies indicate recent or acute infection: HAV IgM, HBV core IgM, VCA IgM (EBV), CMV IgM, measles/mumps IgM, parvovirus B19 IgM, WNV IgM, Zika IgM, dengue IgM, Lyme IgM (with IgG immunoblot confirmation for early disease <30 days of symptoms), \textit{Mycoplasma pneumoniae} IgM
  \item \textbf{Diagnosis of past or chronic infection:} Syphilis (treponemal IgG remains positive for life), HIV (4th-generation antigen/antibody; confirm with differentiation assay + NAT), HCV (anti-HCV with reflex HCV RNA), HTLV-I/II, \textit{Helicobacter pylori} IgG, \textit{Toxocara} IgG, \textit{Strongyloides} IgG
  \item \textbf{Seroconversion testing (paired acute and convalescent sera):} A fourfold or greater rise in specific antibody titer is diagnostic. Used for: \textit{Legionella} (IFA), \textit{Chlamydia psittaci} (CF), \textit{Coxiella burnetii} (Q fever, IFA phase I and II), \textit{Bartonella henselae} (cat scratch disease, IFA), \textit{Mycoplasma pneumoniae} (CF or EIA), arboviruses (WNV, SLE, EEE, WEE, La Crosse), \textit{Coccidioides} (CF, immunodiffusion), \textit{Histoplasma} (CF, immunodiffusion)
  \item \textbf{Autoantibody testing:} ANA (SLE, systemic sclerosis, mixed connective tissue disease, Sjögren), RF/anti-CCP (rheumatoid arthritis), ANCA (granulomatosis with polyangiitis, microscopic polyangiitis, eosinophilic GPA), anti-dsDNA, anti-Sm, anti-RNP, anti-Ro/SSA, anti-La/SSB, anti-Scl-70, anti-centromere, anti-Jo-1, and many others
\end{itemize}

\section*{Primary vs Secondary Test}
Serology is the primary test for determining immunity status, for diagnosis of many infections where PCR or culture is unavailable or impractical (syphilis, Lyme disease, EBV, CMV, leptospirosis, rickettsial diseases, endemic mycoses), and for diagnosis of autoimmune diseases. Serology is a secondary/adjudicative test when PCR or culture is the primary test but is negative or unavailable (paired serology for arboviruses, \textit{Coxiella}, \textit{Bartonella}). For HIV, the 4th-generation combination antigen/antibody assay is the primary screening test.

\section*{How This Test Is Performed}
Venous blood is collected in an SST (serum separator, gold-top or red-top) tube. 5--7 mL of whole blood yields 2--3 mL of serum, sufficient for most serologic panels. The specimen should be allowed to clot for 30 minutes, centrifuged, and serum separated within 2 hours of collection. Serum can be refrigerated at 2--8\textdegree{}C for up to 7 days or frozen at -20\textdegree{}C to -80\textdegree{}C for long-term storage. For paired serology (acute and convalescent), the acute serum should be collected as early as possible in the illness (first week) and the convalescent serum 2--4 weeks later. Both specimens should be tested simultaneously in the same laboratory to avoid inter-assay/inter-laboratory variability. CSF, pleural fluid, and other body fluids can be tested for antibodies when local (intrathecal) antibody production is suspected (CSF IgG index, CSF VDRL, CSF anti-VZV IgG ratio).

\section*{What Preparation Is Needed}
No patient preparation (fasting not required). Recent immunoglobulin therapy (IVIG, VariZIG, HBIG, rabies immune globulin) or blood transfusion can cause passive transfer of antibodies and lead to false-positive serology. A history of such treatments should be documented. The timing of specimen collection relative to symptom onset or exposure is critical for test interpretation; the requesting clinician should provide this information to the laboratory.

\section*{Contraindications and Precautions}
The most important limitation of serology is that it diagnoses the host immune response, not the pathogen itself. Antibodies take time to develop: IgM typically appears 3--7 days after symptom onset for most acute infections, IgG at 7--14 days. Serology collected too early in the illness (within the first 2--3 days) may be falsely negative, even when the patient has the disease. Paired acute and convalescent serology overcomes this limitation. False-positive IgM results occur with: rheumatoid factor (IgM rheumatoid factor can bind to IgG in the assay, producing a false-positive signal); cross-reactivity with other related pathogens (flaviviruses cross-react extensively---Zika, dengue, WNV, YF, JE, SLE, TBE; a positive Zika or dengue IgM in a patient with prior flavivirus infection requires careful interpretation and PRNT for definitive diagnosis); polyclonal B-cell activation during acute EBV, CMV, or HIV infection causing false-positive IgM for many pathogens; and persistence of IgM for months after acute infection (HAV IgM persists 3--6 months; EBV VCA IgM can persist for months; CMV IgM can persist >1 year). False-negative serology can occur in immunocompromised patients (HIV/AIDS with low CD4, transplant recipients, hematologic malignancy, chronic kidney disease, malnutrition) who may be unable to mount a detectable antibody response. In these patients, PCR or antigen detection should be used instead of serology when possible. False-negative serology also occurs in the window period before seroconversion (HIV: 15--30 days; HCV: 6--8 weeks; HBV: up to 60 days for HBsAg to become positive). Serology should NOT be used to monitor treatment response for most infections because antibodies persist long after the pathogen has been cleared. Exceptions: RPR titer for syphilis (declines with successful treatment); HBV anti-HBs titer (monitors vaccine response); HCV antibody reversion (rare, but can occur after successful treatment or in immunosuppressed patients who clear the virus).

\section*{Risks}
Standard venipuncture risks. No risk from the serologic tests. The largest risk is clinical misinterpretation of serology results leading to inappropriate diagnosis and therapy.

\section*{What Happens After the Test}
Serology results are typically available in 1--24 hours (routine assays), with some specialized send-out tests (PRNT, immunoblot, paired serology) requiring 1--2 weeks. Results should be interpreted by a clinician familiar with the patient's clinical presentation, timing of illness, epidemiologic exposures, vaccination history, and immune status.

\section*{Understanding Results}

\subsection*{Below Normal}
\begin{itemize}
  \item \textbf{Negative (nonreactive):} No specific antibodies detected. For immune status screening (VZV IgG negative, measles IgG negative, HAV IgG negative, HBsAb <10 mIU/mL): indicates susceptibility to infection. Vaccination or other preventive measures should be offered. For acute infection diagnosis: a negative result may indicate: absence of infection; specimen collected too early (before seroconversion); immunosuppressed patient unable to mount antibody response; or the wrong serologic panel was ordered
  \item \textbf{Seroconversion:} Negative acute serum with positive convalescent serum indicates recent infection. Diagnostic (e.g., HIV seroconversion, Lyme seroconversion)
\end{itemize}

\subsection*{Above Normal}
\begin{itemize}
  \item \textbf{Positive (reactive) IgM:} Indicates recent or acute infection. Interpretation caveats: (a) IgM may persist for months and does not necessarily indicate infection within the past few days; (b) IgM can be falsely positive (rheumatoid factor, cross-reactivity); (c) IgM in a newborn indicates congenital infection (IgM does not cross the placenta). For infections where IgM is diagnostic (HAV, HBV core, EBV VCA, parvovirus B19, measles, rubella), further testing is typically not required unless the clinical picture is atypical
  \item \textbf{Positive (reactive) IgG:} Indicates past infection or vaccination. Depends on the pathogen: VZV IgG positive = immune to varicella; measles IgG positive = immune; CMV IgG positive = past infection (latent, with risk of reactivation if immunosuppressed); HCV IgG positive = past or current HCV; requires HCV RNA to distinguish. A positive IgG does not distinguish recent from remote infection unless paired with acute serum demonstrating seroconversion or a fourfold titer rise. IgG avidity testing (low avidity = recent <3--6 months; high avidity = remote) can help when this distinction is clinically important (CMV in pregnancy, toxoplasmosis in pregnancy)
  \item \textbf{Fourfold or greater rise in titer:} The gold standard for serologic diagnosis of recent infection. Requires paired acute and convalescent sera tested simultaneously. Used when direct detection (PCR, culture) is unavailable or negative. Examples: rickettsial infections, \textit{Coxiella burnetii}, \textit{Bartonella}, \textit{Mycoplasma pneumoniae}, \textit{Chlamydia psittaci}, arboviruses, endemic mycoses (CF titers for \textit{Histoplasma}, \textit{Coccidioides})
  \item \textbf{Indeterminate/equivocal:} The result is near the cutoff between negative and positive. May represent: early seroconversion (repeat in 2--4 weeks); waning antibody levels; nonspecific cross-reactivity; or a technical issue. Repeat testing in 2--4 weeks or testing by an alternative method is recommended
\end{itemize}

\section*{Normal Results}
\textbf{Nonreactive (Negative)} or specific protective antibody level: anti-HBs $\geq$10 mIU/mL (immune to HBV). For autoantibodies: \textbf{Negative} (titer below the laboratory's reference range). For immune status: depends on the pathogen---positive IgG (immune) or negative IgG (susceptible).

\section*{Follow-up Tests}
\begin{itemize}
  \item \textbf{Convalescent serology (repeat testing at 2--4 weeks):} When acute serology is negative, for seroconversion or a fourfold titer rise. Essential for many vector-borne and zoonotic infections
  \item \textbf{IgG avidity testing:} When the timing of infection is critical (CMV, toxoplasmosis in pregnancy). Low avidity = primary infection within the preceding 3--6 months
  \item \textbf{Confirmatory/alternative serologic method:} When the initial result is equivocal or discordant with the clinical picture (e.g., confirm low-positive HSV-2 IgG with HSV-2 immunoblot; confirm positive HIV screen with differentiation assay and NAT; confirm positive syphilis treponemal EIA with TPPA)
  \item \textbf{PCR (nucleic acid amplification testing):} When serology is negative but clinical suspicion is high and the infection is in the early window period (HSV encephalitis, acute HIV, early HCV, early Lyme with erythema migrans), PCR should be performed directly on the affected site
  \item \textbf{Neutralization assays (PRNT, HAI):} Gold standard to confirm specificity of serologic results when cross-reactivity is a concern (flavivirus serology) or to measure functional, protective antibody levels
\end{itemize}

\section*{References}
\bibliographystyle{unsrtnat}
\bibliography{references}

\clearpage
